% !TEX program = xelatex
\documentclass{英语试卷}
\begin{document}
\examtitle{泉州市 2025 届高中毕业班质量监测(四)}{英语}
\par 泉州市 2025 届高中毕业班质量监测(四)\par
\par 高三英语\par
\begin{notice}
\par 1. 答卷前，考生务必将自己的姓名、准考证号填写在答题卡上。\par
\par 2. 回答选择题时，选出每小题答案后，用铅笔把答题卡上对应题目的答案标号涂黑。如需改动，用橡皮擦干净后，再选涂其他答案标号。回答非选择题时，将答案写在答题卡上，写在本试卷上无效。\par
\par 3. 考试结束后，将本试卷和答题卡一并交回。\par
\end{notice}
\parttitle{第一部分 听力（共两节，满分 30 分）}
\par 做题时，先将答案标在试卷上。录音内容结束后，你将有两分钟的时间将试卷上的答案转涂到答题卡上。\par
\sectitle{第一节（共5小题：每小题1.5分，满分7.5分）}
\par 听下面 5 段对话。每段对话后有一个小题，从题中所给的 A、B、C 三个选项中选出最佳选项。听完每段对话后，你都有 10 秒钟的时间来回答有关小题和阅读下一小题。每段对话仅读一遍，\par
\question{1. What does Mr. Thompson do for the community?}
\begin{choices}(3)
  \item He attends to patients in the hospital.
  \item He helps to educate children.
  \item He donates to the community center.
\end{choices}
\question{2. What are the speakers mainly talking about?}
\begin{choices}(3)
  \item A meeting.
  \item A new technique.
  \item A place for healthcare.
\end{choices}
\question{3. What does the woman want to be?}
\begin{choices}(3)
  \item A film actress.
  \item A stage designer.
  \item A theater director.
\end{choices}
\question{4. What are the speakers doing?}
\begin{choices}(3)
  \item Gardening.
  \item Shopping.
  \item Cooking.
\end{choices}
\question{5. To whom does the woman give credit?}
\begin{choices}(3)
  \item Her class.
  \item Her team.
  \item Her teacher.
\end{choices}
\sectitle{第二节（共15小题：每小题1.5分，满分22.5分）}
\par 听下面 5 段对话或独白。每段对话或独白后有几个小题，从题中所给的 A、B、C 三个选项中选出最佳选项。听每段对话或独白前，你将有时间阅读各个小题，每小题 5 秒钟；听完后，各小题将给出 5 秒钟的作答时间。每段对话或独白读两遍。 听第 6 段材料，回答第 6、7 题。\par
\question{6. What does the woman love most about Ne Zha?}
\begin{choices}(3)
  \item The lines.
  \item The images.
  \item The theme.
\end{choices}
\question{7. What contributes to Ne Zha's success?}
\begin{choices}(3)
  \item Designing around target audience.
  \item Borrowing ideas from ancient stories.
  \item Vitalizing traditional tales. 听第 7 段材料，回答第 8 至 10 题。
\end{choices}
\question{8. What does the woman see as the key to positive relationships?}
\begin{choices}(3)
  \item Wearing a friendly smile.
  \item Approaching others willingly.
  \item Being a good listener.
\end{choices}
\question{9. What does the woman recommend to do during conversations?}
\begin{choices}(3)
  \item Impress people intentionally.
  \item Maintain eye contact.
  \item Avoid disagreements.
\end{choices}
\question{10. How does the man respond to the advice?}
\begin{choices}(3)
  \item He approves it.
  \item He doubts it.
  \item He dismisses it.
\end{choices}
\par 听第 8 段材料，回答第 11 至 13 题，\par
\question{11. What time is the last bus on weekdays?}
\begin{choices}(3)
  \item 6 p.m.
  \item 7 p.m.
  \item 8 p.m.
\end{choices}
\question{12. What is the probable relationship between the speakers?}
\begin{choices}(3)
  \item Strangers.
  \item Fellow classmates.
  \item Driver and passenger.
\end{choices}
\question{13. Which pickup point is more convenient for the woman?}
\begin{choices}(3)
  \item The library.
  \item The health clinic.
  \item The Happy Supermarket.
\end{choices}
\par 听第 9 段材料，回答第 14 至 17 题。\par
\question{14. What is the topic of the conversation?}
\begin{choices}(3)
  \item Travelling in California.
  \item Adaptation to college life.
  \item Ways of pressure management.
\end{choices}
\question{15. What challenges Sophia most on the debate team?}
\begin{choices}(3)
  \item Preparing for debates.
  \item Fitting in with the team.
  \item Managing time effectively.
\end{choices}
\question{16. What makes the debates more interesting?}
\begin{choices}(3)
  \item The members' close relationship.
  \item The members' diverse viewpoints.
  \item The members' debating techniques.
\end{choices}
\question{17. How does Sophia achieve balance?}
\begin{choices}(3)
  \item By keeping to a timetable.
  \item By relaxing herself regularly.
  \item By following friends' track.
\end{choices}
\par 听第 10 段材料，回答第 18 至 20 题。\par
\question{18. What is the scientific method used for?}
\begin{choices}(3)
  \item Researching the natural world.
  \item Confirming personal opinions.
  \item Inventing new technologies.
\end{choices}
\question{19. What example is given for observation?}
\begin{choices}(3)
  \item The duration of sunlight in a certain area.
  \item The rate at which temperature changes.
  \item The effect of sunlight on temperature.
\end{choices}
\question{20. What is the purpose of the talk?}
\begin{choices}(3)
  \item To argue.
  \item To inform.
  \item To advertise.
\end{choices}
\parttitle{第二部分 阅读理解（共两节，满分 50 分）}
\sectitle{第一节（共15小题：每小题2.5分，满分37.5分）}
\par 阅读下列短文，从每题所给的 A、B、C 和 D 四个选项中，选出最佳答案。\par
\passagelabel{A}
\par The 2025 Metropolis Marathon will take place on December 1 in the districts of Downtown, Riverside, and Uptown. For interested runners, the following is a guide for registration.\par
\par Registration Period\par
\par From October 9 at 8 pm to October 14 at noon, participants can log on to either the official website of the Metropolis Marathon or the MetMar app. They can then select the corresponding races and events for registration. A drawing will be held, for which the results will be announced on the afternoon of October 24. After the results are announced, participants have to complete the registration payment before October 28 at noon.\par
\par Registration Fee\par
\par For foreign runners, the registration fee for the full marathon is 235 per person, 145 per person for the half marathon, and 80 for the fitness running program. Domestic runners pay 225, 135, and 70 respectively.\par
\par Entry Requirements\par
\par Direct entry is offered to top finishers of previous Metropolis races (e. g., Women's Half Marathon, Riverside Half Marathon). Additionally, 100 participants from the 2024 Virtual Run can enter via lottery(抽签). Qualified runners must register between October 9 and October 17.\par
\par Charity Program\par
\par This year, the Metropolis Marathon will continue to offer charity spots. Registration opens at 8 pm on October 9 with a total of 800 seats available on a first-come, first-served basis. Charity runners who successfully register are automatically confirmed and do not need to enter the drawing.\par
\par Key Notes\par
\par Late payments will result in disqualification.\par
\par International runners must provide valid(有效的) passports and medical reports.\par
\par No refunds are allowed after October 28.\par
\question{21. When will the registration result be released?}
\begin{choices}(2)
  \item On Oct 9.
  \item On Oct 14.
  \item On Oct 17.
  \item On Oct 24.
\end{choices}
\question{22. Who can get direct entry?}
\begin{choices}(2)
  \item Athletes of the 2024 Virtual Run.
  \item Participants of fitness running program.
  \item Winners of past Metropolis races.
  \item Foreign runners paying \$235.
\end{choices}
\question{23. What can a runner do to secure a charity spot?}
\begin{choices}(2)
  \item Enter the runner lottery.
  \item Register within 800 seats.
  \item Pay an additional fee.
  \item Submit a charity application.
\end{choices}
\passagelabel{B}
\par In Baltimore, Our Kids Read, a nonprofit offers three free books to every child who walks into its bookstore. And unlike the library, the books are never expected back. Kids "check out" at the front of the store, but instead of handing over cash, all they have to do is sign up for the nonprofit’s Reading Buddy program, which matches young readers with adults who read with them once or twice a week in 45-minute online meetings to promote literacy (读写能力). "Every child should have someone reading to them at night," said Jahmal Lake, executive director of Our Kids Read.\par
\par Lake established Our Kids Read in 2021. But he quickly realized that he needed to do more than just distribute donated books. "I think that the mistake a lot of nonprofits make, is that they think they can back up a truck full of resources to the communities they want to serve and that they're doing something," Lake said. "But if that community is not equipped to make good use of those resources, you might as well have given them a brick."\par
\par On the bookshelves and across tables, there are books like "Black Boy Joy" and "Eyes That Kiss in the Corners", picture books in English and Spanish, poetry collections, and even graphic novels. The range of options plays into Lake's vision of turning non-readers into bookworms by directing kids to books on subjects they already love. Research shows kids are more likely to finish and love books they pick out themselves.\par
\par To promote the Reading Buddy program, Our Kids Read goes into elementary schools. Lake also plans to set up several little free libraries to encourage reading. The Baltimore Community Foundation has given Our Kids Read over \$25,000. Vice President of Community Impact Crystal Harden-Lindsey said Our Kids Read is vital to promoting literacy, the foundation for academic success, economic mobility and civic engagement.\par
\question{24. What is a feature of Our Kids Read bookstore?}
\begin{choices}(2)
  \item It offers free books to adults and kids.
  \item It holds reading competitions weekly.
  \item It engages kids in reading books to adults.
  \item It requires kids to register for a program.
\end{choices}
\question{25. What does Lake think of merely donating books?}
\begin{choices}(2)
  \item Insufficient.
  \item Meaningful.
  \item Motivating.
  \item Unnecessary.
\end{choices}
\question{26. Why does the bookstore offer diverse books?}
\begin{choices}(2)
  \item To expand kids' reading range
  \item To develop students' love for books
  \item To broaden kids' visions on book selection.
  \item To encourage kids' interest in school subjects.
\end{choices}
\question{27. What is the main idea of the passage?}
\begin{choices}(2)
  \item Bookstore offers kids free books to promote business.
  \item Charity partners with schools to teach students reading.
  \item Nonprofit boosts child literacy through multiple measures.
  \item Program enhances family bond by shared reading activities.
\end{choices}
\passagelabel{C}
\par Early solar expansion often came at agriculture's expense. Farmers and ranchers watched as fields once used for grazing(放牧) were fenced off, their livelihoods threatened by the growth of renewable energy. By the 2010s, conflicts peaked: in Texas alone, over 200,000 acres of grassland were changed into solar farms, reducing grazing capacity by nearly 40\%, generating debates over "clean energy vs cattle."\par
\par Then came a transformative idea: agrivoltaics. Pioneered in the 2020s, this approach reimagined solar farms not as competitors to agriculture, but as partners. The concept combines solar energy production with sustainable land use. JR Howard, a Texas rancher who first offered his sheep to solar operators in 2021, later founded Texas Solar Sheep—a booming business tending over 8,000sheep across multiple solar sites, employing dozens of workers in rural communities where jobs were rare. His story reflects a broader restoration, as agricultural traditions find new purpose in the renewable energy era.\par
\par The benefits extend beyond economics. Each grazing sheep represents one less gallon of gasoline used by mowers(割草机), a part of habitat preserved for ground-nesting birds, and a living proof of sustainable land management. Researchers note another unexpected advantage - the solar panels provide shade that improves animal welfare, creating microclimates where sheep grow healthily even in Texas' heat.\par
\par This marriage of agriculture and energy now covers over 60 projects nationwide. The U.S. Department of Agriculture reports a 22\% rise in lamb production in solar-grazing regions since 2020, alongside a 35\% reduction in land maintenance costs for energy companies.\par
\par However, important questions remain. Scientists like Nuria Gomez-Casanovas from Texas A\&M stress that the long-term effects on soil health and food production need more study. Will these double-use lands maintain fertility (肥沃)? Will this model be effective in different climates and agricultural systems?\par
\par Nonetheless, agrivoltaics shows a promising model of coexistence. Each solar panel generates energy and serves as shelter for farm animals during heatwaves. Sheep droppings improve soil, enhancing grassland health. As solar fields and farm animals grow well together, agrivoltaics proves that progress doesn't have to erase tradition—it can support both.\par
\question{28. What initially caused conflicts between solar farms and farmers?}
\begin{choices}(2)
  \item The abandonment of farmland.
  \item Pollution from solar panel production.
  \item The high cost of gas-powered mowers.
  \item The occupation of grazing land by solar panels.
\end{choices}
\question{29. What does solar grazing bring to farmers?}
\begin{choices}(2)
  \item New income opportunities.
  \item Increased lamb prices.
  \item Higher land maintenance costs.
  \item Free solar energy access.
\end{choices}
\question{30. What can be inferred about agrivoltaics' environmental impact?}
\begin{choices}(2)
  \item It relies on fuel-driven equipment.
  \item It balances energy and ecology.
  \item It harms soil quality by overgrazing.
  \item It increases competition for resources.
\end{choices}
\question{31. Which title best summarizes the passage?}
\begin{choices}(2)
  \item The Rise of Solar Energy in Texas
  \item Challenges in Grazing Industry
  \item Agrivoltaics: Where Sheep and Solar Panels Meet
  \item Agrivoltaics: How Tech Innovation Connects Farms
\end{choices}
\passagelabel{D}
\par Artificial intelligence has persuaded a lot of folks that we need to undertake a comprehensive reform of education. Now that chatbots can speedily access information and answer complex questions, why bother memorizing historical facts or quotations? Shouldn't we instead teach kids to think critically and solve problems and leave the boring duties to computers?\par
\par There's a problem with these arguments: humans require knowledge to think creatively. Outsourcing our memory and cognition(认知) to Google and AI risks making us less intelligent and more subject to false information.\par
\par In 2000, American educator E. D. Hirsch noted in his essay: "The internet places information at our fingertips, but using it demands existing knowledge stores." This principle drove the UK's 2014 curriculum reform empowering multiplication table mastery(乘法口诀表) by age nine-not merely for calculation skills, but to build the neural frameworks (神经框架) enabling complex reasoning. As former schools minister Nick Gibb stressed: “‘Google it effectively depends on what’s stored in long-term memory."\par
\par "The suggestion we should outsource our memory to "free up" limited space for more creative thinking is based on a misconception—we're likely to overestimate our intelligence, mistaking knowledge found on the internet for our own," writes Nicholas Carr in The Shallows. "When we start using the Web as a replacement for personal memory, bypassing the inner processes of consolidation(巩固), we risk emptying our minds of their riches."\par
\par Indeed, looking stuff up on Google often results in us not being able to recall it later—either because our brains are conditioned to think we don't need to remember it, or because the internet, mobile phones and social media distract our attention, or both. If students come to rely on machines not just to recall facts but to think for them, it might not just be memory that suffers: cognition and creativity could atrophy, too.\par
\par The solution lies not in rejecting technology, but in rebalancing its role. The UK's knowledge-rich curriculum, despite disagreements, recognizes this truth: human intelligence advances when technology adds to rather than replaces learning.\par
\question{32. What is the author's concern about AI's influence on education?}
\begin{choices}(2)
  \item Students avoid complex questions.
  \item Teachers prioritize memorizing facts.
  \item Students develop poor analysis and judgment.
  \item Teachers ignore students' creativity.
\end{choices}
\question{33. What do experts consistently emphasize?}
\begin{choices}(2)
  \item The necessity for technology improvement.
  \item The importance of knowledge storage.
  \item The reduction of internet usage.
  \item The replacement of complex reasoning.
\end{choices}
\question{34. What does the underlined word "atrophy" in paragraph 5 mean?}
\begin{choices}(2)
  \item Decline.
  \item Disappear.
  \item Strengthen.
  \item Stabilize.
\end{choices}
\question{35. Which of the following may the author agree?}
\begin{choices}(2)
  \item Limit AI to education.
  \item Avoid referring to AI in study.
  \item Search online for solutions beforehand.
  \item Use technology as an assistant.
\end{choices}
\sectitle{第二节（共5小题；每小题2.5分，满分12.5分）}
\par 阅读下面短文，从短文后的选项中选出能填入空白处的最佳选项。选项中有两项为多余选项，\par
\par If you want to perform at a high level, you'll need to be resilient (韧性). The way to build this mental skill is by overcoming unexpected challenges that come your way. \blanknum{36} High performers approach challenges differently, often inviting them in and even using mental frameworks that fundamentally change the way they see the world.\par
\par Do Something Hard Every Day\par
\par Extreme athletes are some of the most resilient individuals among us. \blanknum{37} For example, top athletes often use elective hardship in their lives. In this way, they are better prepared for the tough moments in life that aren't voluntary.\par
\par Look Pain in the Face\par
\par \blanknum{38} Courtney Dauwalter, an endurance athlete, thinks about pain as an actual place she calls a "pain cave". She knows whenever she runs really long races, she'll eventually reach the pain cave. "If I pull through the darkness, I leave as a stronger self," she says. It's a process of transformation.\par
\par Redefine Hardships\par
\par It also helps to put moments of hardship in a positive light. Tommy Caldwell, a rock climber, looks at moments of pain as an adventure. \blanknum{39} "If you can reframe these moments, then you're better able to deal with them and pull through," he said. For him, physical challenges help him feel good in his body, which then helps him feel good in his mind.\par
\par Mental toughness is a skill. It's like a muscle that you exercise for the coming terrible moments. If you introduce voluntarily challenges into your everyday life, you're not as fragile (脆弱的) when those big moments come. \blanknum{40}\par
\begin{choices}(1)
  \item Their methods model resilience-building strategies.
  \item Unexpectedly as they may come, they also go easily.
  \item Another way to build resilience is by personifying pain.
  \item They have the ability to overcome pain for long periods of time.
  \item This mindset shift enables him to tackle barriers with confidence.
  \item And you can rely on the fact that you have the skills to face that moment.
  \item While we all experience roadblocks, you don't have to wait for one to happen.
\end{choices}
\parttitle{第三部分 语言知识运用（共两节，满分 30 分）}
\sectitle{第一节（共15小题；每小题1分，满分15分）}
\par 阅读下面短文，从短文后各题所给的 A、B、C 和 D 四个选项中，选出可以填入空白处的最佳选项，\par
\par From Canada comes the story of a lost artwork, a bargain-bin (廉价品区) treasure of a lifetime and a connection that grew into a friendship across 1,200 miles.\par
\par Having spent 300 hours each \blanknum{} 41 \blanknum{} two plates of color-phase tiles(变色瓷砖), Nova Scotian artist Sydney Blum was excited at the chance to \blanknum{} 42 \blanknum{} her work in a Montreal gallery. She shipped them, but both got \blanknum{} 43 \blanknum{} in the mail. Canada Post failed to locate the package, leaving Blum \blanknum{} 44\par
\par Weeks later, Sonja Krawesky, a teacher, was \blanknum{45} by two colorful plates when she was looking through items at a bargain-bin store in Hamilton. Krawesky loved their \blanknum{46}. She knew these did not \blanknum{47} there and took them home to find out their \blanknum{48}.\par
\par One morning, an email about the lost artwork reached Blum. "I lost my 49 when I read it," Blum said.\par
\par Krawesky had \blanknum{50} Blum through the Montreal gallery, and discovered the \blanknum{51} of both the plates and Sydney Blum.\par
\par Blum \blanknum{52} she was being scammed(诈骗) at first. After confirming Krawesky's \blanknum{53} through her school, Blum was blown away. But Krawesky suggested not sending the plates by mail again. \blanknum{54}, Blum contacted a trucker to take them to Nova Scotia.\par
\par The encounter developed into a \blanknum{55}. "It is special." Krawesky said. "I am motivated to do more of my own creative things."\par
\cz{41}{creating}{assessing}{decorating}{packaging}
\cz{42}{repair}{exhibit}{sell}{design}
\cz{43}{broken}{mixed}{burned}{lost}
\cz{44}{guilty}{hesitant}{heartbroken}{merciful}
\cz{45}{attracted}{bothered}{touched}{stimulated}
\cz{46}{usage}{price}{company}{look}
\cz{47}{exist}{work}{belong}{match}
\cz{48}{destination}{origin}{history}{value}
\cz{49}{way}{memory}{patience}{breath}
\cz{50}{tracked down}{referred to}{came across}{waited for}
\cz{51}{secret}{information}{nature}{story}
\cz{52}{admitted}{realized}{worried}{explained}
\cz{53}{intention}{significance}{background}{identity}
\cz{54}{However}{Instead}{Meanwhile}{Moreover}
\cz{55}{culture}{clue}{bond}{dream}
\sectitle{第二节（共10小题；每小题1.5分，满分15分）}
\par 阅读下面短文，在空白处填入 1 个适当的单词或括号内单词的正确形式。\par
\par Recently, a video has gone viral online. On the busy streets of Sydney, Australia, a young Chinese woman, Dong Min, dressed in beautiful hanfu, the traditional clothing of China, elegantly \blanknum{56} (hold) a dizi, the bamboo flute of China, inviting passers-by with traditional Chinese music.\par
\par Born in 1990, Dong graduated from Minzu University of China with \blanknum{57} degree in the major of Chinese bamboo flute and later studied at the University of New South Wales in Australia, \blanknum{58} she formed a Chinese music band with like-minded friends. \blanknum{59} (subsequent), she founded the Sydney Meya Conservatory of Chinese Music, \blanknum{60} (devote) to teaching Chinese traditional instruments to more people. It attracts students aged 4 to 80, reflecting the timeless appeal of Chinese music across; \blanknum{61} (generation). Having noticed a lack of awareness of Chinese traditional music in Australia, she \blanknum{62} (inspire) by local performers using western instruments like pianos and violins \blanknum{63} (launch) street performances with her members, showing traditional Chinese instruments in Sydney's Chinatown in 2016.\par
\par "Street performances offer a direct and visual way for people to experience Chinese music, which makes it \blanknum{64} (access) to everyone," she said. As Dong's performances grow in popularity, videos of her shows spread across social media, drawing in a larger audience. Dong has witnessed a change in foreign audiences, from unfamiliarity \blanknum{65} active engagement, demonstrating the growing presence of traditional Chinese music in Australia, both in formal settings and on the streets.\par
\parttitle{第四部分 写作（共两节，满分 40 分）}
\sectitle{第一节（满分15分）}
\par 假定你是李华，上周你校举办了“劳模校友进校园”宣讲活动。请你以"Learning from Model Workers"为题，给你校英文报 Youth 写一篇短文投稿，内容包括：(1) 劳模故事；(2) 劳模精神；(3) 感受和启发。注意：\par
\par (1) 写作词数应为 80 左右;\par
\par (2) 请按如下格式在答题纸的相应位置作答。\par
\par Learning from Model Workers\par
\sectitle{第二节（满分25分）}
\par 阅读下面材料，根据其内容和所给段落开头语续写两段，使之构成一篇完整的短文，\par
\par The sweet smell of lemon cleaning spray mixed with the light smell of oil paints as Lila and her two classmates pushed open the heavy doors of Cedar Hill Nursing Home. Their art teacher, Ms. Reed, had arranged this volunteer program — three Saturdays spent helping redecorate the activity room.\par
\par "Careful with that ladder, Raj!" Lila warned as her tech-loving friend nearly tripped over a pot of brushes. Maya, always the quiet observer, stopped to draw the sunlight shining through the dusty curtains.\par
\par As they pushed an old bookshelf to paint the wall behind it, Maya suddenly froze. "Look behind this shelf! These aren't ordinary marks..."\par
\par The three teens pulled the wooden shelf aside, dust dancing in the golden afternoon light. Faded across the wall were beautiful pencil drawings: a teapot growing cherry blossoms (樱花), a bicycle flying above clouds, and a woman dancing with twin girls under a pink sky.\par
\par Raj took photos with his phone. "The lines are shaky-like drawn in a hurry."\par
\par Nurse Clara rushed over. "Oh dear, you found Mrs. Talbot's secret drawings. Whenever words fail her, she draws memories at night." She sighed. "We've seen marks before, but never this many."\par
\par Lila's paintbrush froze above the white paint can. Their deadline was near, yet covering these lines felt like silencing a singing bird. Her throat tightened as Maya grabbed her arm. "We can't cover this! It's like... like erasing her diary!"\par
\par Raj adjusted his glasses. "What if we scan (扫描) them? My 3D printer could make models for the elders to touch!"\par
\par "We must save these!" Maya's voice trembled as she reached toward the wall. Her fingertips brushed the cherry blossoms. "But how? What if we keep them... in a new way?"\par
\par Ms. Reed joined them, listening intently. "We'll need permission," she said, "but this could be special."\par
\par Lila wiped her hands on her jeans and suggested, "I'll go with Nurse Clara to talk to the manager."\par
\par Raj proposed scanning the pictures with phones, while Maya suggested adding animated (动画的)falling flowers that could play voices. Ms. Reed offered to use colored lights for projection(投影). 注意：\par
\par (1) 续写词数应为 150 左右;\par
\par (2) 请按如下格式在答题卡的相应位置作答。\par
\writingpara{Paragraph 1: Over the next week, they worked tirelessly.}
\writinglines{8}
\writingpara{Paragraph 2: On that golden Saturday, a digital exhibition lit the room with living memories.}
\writinglines{8}
\end{document}
